% Solutions to Chapter 5's exercises, from lectures/L05/appendix/c_solutions.md. No exercise in the
% chapter is a program, so every one is answered, and the source's note on what counts as the
% answer in a troubleshooting case is kept as the entry's introduction.

\answerchapter{c5:ch}

I felsökningsfallen är det \textbf{resonemanget} som är svaret: vilken mätning som görs härnäst, och
vad den kan visa. Den som gissade rätt orsak utan att kunna säga var hen skulle mäta har gissat.

\answer{c5:ex:1}{Matningen först}
\pt{a}Röd mätsladd på \textbf{ben~14} (VCC), svart på \textbf{ben~7} (GND). Du ska se ungefär
\textbf{5~V}.

\pt{b}Skenan kan ha 5~V medan kretsen inte har det: sladden till ben~14 kan sitta i kolumnen
bredvid, eller skenan kan vara delad på mitten av kopplingsdäcket. Det som räknas är spänningen
\textbf{på kretsens ben}, så det är där du mäter.

\pt{c}Varje ingång på en CMOS-krets har skyddsdioder mot VCC. När en ingång får 5~V från en
strömbrytare leder dioden, och kretsen får en del av sin matning \textbf{genom ingången}. Kretsen
kan då ge utsignaler som ibland är rätt, men bara när någon ingång råkar vara hög, och med fel
nivåer. Det är därför ett matningsfel sällan ser ut som \enquote{ingenting fungerar}, och därför
matningen ska mätas först även när något verkar fungera.

\answer{c5:ex:2}{Vad betyder mätvärdet?}
\pt{a}4,7~V är över 3,5~V: logisk \textbf{1}, en normal hög utgång.

\pt{b}0,1~V är under 1,5~V: logisk \textbf{0}, en normal låg utgång.

\pt{c}2,3~V ligger i det förbjudna området mellan 1,5 och 3,5~V: \textbf{ingen giltig nivå}.
Troligen en flytande ingång, som inte är ansluten till något.

\pt{d}2,5~V på en utgång som är kopplad till en annan utgång: \textbf{två utgångar mot varandra}.
Den ena försöker driva hög, den andra låg, och spänningen hamnar någonstans mitt emellan medan en
onödigt stor ström går mellan dem. Det är ett fel, och kretsarna blir varma. Två utgångar får
aldrig kopplas ihop.

\answer{c5:ex:3}{En grind som ljuger}
\pt{a}En AND-grind med ingångarna 1 och 0 ska ge \textbf{0}, alltså under 0,3~V på ben~3.

\pt{b}Grinden ger 1 för ingångarna 1 och 0: den beter sig som en \textbf{OR-grind}. Den troligaste
förklaringen är att det sitter en \textbf{74HC32} i stället för en 74HC08. De har samma
benplacering, så nätet går att koppla upp och fungerar på alla rader där AND och OR råkar ge samma
svar (00 och 11), vilket gör felet lätt att missa. Kontrollera texten på kretsen. Ett alternativ är
att ben~3 är kortslutet mot 5~V; det skiljer du ut genom att ställa ingångarna på 00: en OR ger då
0, en kortslutning ger fortfarande 5~V.

\answer{c5:ex:4}{Den varma kretsen}
\pt{a}\textbf{Bryt matningen direkt.} En varm krets håller på att skadas, och det går fort.

\pt{b}De två vanligaste orsakerna:
\begin{itemize}
  \item \textbf{kretsen sitter felvänd}, så att den får plus på ben~7 och jord på ben~14;
  \item \textbf{en utgång är kortsluten}: mot matningen, mot jord, mot en annan utgång eller mot en
    strömbrytare som driver den i motsatt riktning.
\end{itemize}

\pt{c}Kontrollera först riktningen: hacket åt vänster, ben~1 längst ned till vänster, och ben~14 i
kolumnen med sladden till \code{+}. Om riktningen är rätt, gå igenom kopplingslistan rad för rad
och leta efter en ledning som går till ett utgångsben (ben 3, 6, 8 och 11 på en 74HC08) från något
annat än en ingång. Multimeterns resistansläge (Ω), med matningen \textbf{frånslagen}, visar om ett
utgångsben har direkt förbindelse med en skena.

\answer{c5:ex:5}{Läs felet ur tabellen}
\pt{a}Förväntat är \code{X = 1} på raderna 001, 011, 101, 110 och 111. Raderna \textbf{001, 011
och 101} är fel: där blir \code{X = 0}.

\pt{b}Alla tre har \code{C = 1} och \code{AB = 0}. Det är de rader där \code{X} bara kan bli 1
\textbf{tack vare C}. Raderna där \code{AB = 1} stämmer. Alltså når \code{AB} fram genom
OR-grinden, men \code{C} gör det inte: den uppmätta tabellen är precis \code{X = AB}.

\pt{c}Ställ in raden \code{ABC = 001} och mät på \textbf{OR-grindens andra ingång}, U2 ben~2.
Förväntat 1. Om den visar 0 sitter felet i ledningen från strömbrytaren C: fel kolumn, eller
kanske kopplad till jord tillsammans med de oanvända ingångarna. Om den visar 1 är det OR-grinden
eller dess utgång som ska undersökas.

\answer{c5:ex:6}{Bromsassistenten som inte bromsar}
\pt{a}Kravet att \textbf{föraren alltid ska kunna bromsa}, även när assistanssystemet är trasigt.
Nätet släpper inte igenom förarens pedal när \code{F = 1}, alltså precis när systemet rapporterar
ett fel. Just då, när föraren inte kan lita på systemet och bromsar själv, bromsar bilen inte. Det
är den farligaste raden i hela tabellen.

\pt{b}Pedalen är kopplad \textbf{in i OR-grinden tillsammans med detektorerna}, så att den passerar
AND-grinden med \code{F'}. Nätet de byggt är:

\begin{textdrawing}
B = (D + S + R) · F'
\end{textdrawing}

\noindent i stället för det rätta:

\begin{textdrawing}
B = D + (S + R) · F'
\end{textdrawing}

\noindent Kontrollera: med \code{D = 1} och \code{F = 1} ger det felaktiga uttrycket
\code{1 · 0 = 0}, och det rätta \code{1 + ... = 1}. Med \code{F = 0} ger båda samma svar, vilket är
varför tolv rader stämmer.

\pt{c}Följ ledningen från strömbrytaren \code{D}. Den ska gå till \textbf{den sista OR-grindens}
ingång, den vars utgång är \code{B}. Ställ in en rad med \code{D = 1, F = 1} och mät på ingångarna
till grinden som driver \code{B}. I det rätta nätet är det OR-grinden, med \code{D = 1} på ena
ingången. I det felaktiga nätet är det AND-grinden med \code{F'}, och dess \code{F'}-ingång ligger
på 0: pedalens etta kommer aldrig förbi den.

\answer{c5:ex:7}{Halvering}
\pt{a}\textbf{Efter grind~8}, i mitten. Då vet du efter en enda mätning i vilken halva felet
sitter, och halva nätet är avskrivet.

\pt{b}Sexton grindar halveras till åtta, fyra, två och en: högst \textbf{fyra mätningar}.

\pt{c}Från ingången och grind för grind: i värsta fall \textbf{15 mätningar} (efter grind 1 till
15; om alla är rätt är det grind~16 som felar).

\answer{c5:ex:8}{Den flimrande lysdioden}
\pt{a}En \textbf{flytande ingång} någonstans i nätet: en ingång som inte är ansluten till något, och
som därför påverkas av laddning från omgivningen, till exempel från din hand.

\pt{b}Mät spänningen på varje ingång i nätet, också på oanvända grindar. En ansluten ingång visar
stadigt nära 0~V eller 5~V; den flytande visar ett värde som vandrar, eller ligger mitt emellan,
och ändras när du rör handen. Jämför sedan med kopplingslistan: vilken ingång saknar en rad?

\pt{c}En HC-krets är byggd i CMOS, och dess ingångar har extremt hög resistans: nästan ingen ström
flyter in eller ut, så ingenting bestämmer spänningen. En LS-krets (TTL) har ingångar där en liten
ström flyter \textbf{ut} ur ingången, och den strömmen drar en okopplad ingång mot hög nivå. En
flytande LS-ingång läses därför nästan alltid som 1. Det gör felet mindre synligt på LS, men det är
fel där också.

\needspace{14\baselineskip}
\answer{c5:ex:9}{Planera fläkten}
\pt{a}Grindtilldelning för \code{F = TW'B' + CB'}:

\begin{narrowtable}{@{}llclc@{}}
  \thead{Grind i nätet} & \thead{Krets} & \thead{Grind} & \thead{Ingångar (ben)} &
    \thead{Utgång (ben)} \\
  \midrule
  NOT: \code{W'} & U1 = 74HC04 & 1 & W $\rightarrow$ 1 & 2 \\
  NOT: \code{B'} & U1 = 74HC04 & 2 & B $\rightarrow$ 3 & 4 \\
  AND: \code{TW'} & U2 = 74HC08 & 1 & T $\rightarrow$ 1, U1 ben 2 $\rightarrow$ 2 & 3 \\
  AND: \code{TW'B'} & U2 = 74HC08 & 2 & U2 ben 3 $\rightarrow$ 4, U1 ben 4 $\rightarrow$ 5 & 6 \\
  AND: \code{CB'} & U2 = 74HC08 & 3 & C $\rightarrow$ 9, U1 ben 4 $\rightarrow$ 10 & 8 \\
  OR: \code{F} & U3 = 74HC32 & 1 & U2 ben 6 $\rightarrow$ 1, U2 ben 8 $\rightarrow$ 2 & 3 \\
\end{narrowtable}

\needspace{20\baselineskip}
\pt{b}Kopplingslista:

\begin{narrowtable}{@{}rll@{}}
  \thead{Nr} & \thead{Från} & \thead{Till} \\
  \midrule
  1 & U1, U2 och U3 ben 14 & +5~V \\
  2 & U1, U2 och U3 ben 7 & GND \\
  3 & strömbrytare W & U1 ben 1 \\
  4 & strömbrytare B & U1 ben 3 \\
  5 & strömbrytare T & U2 ben 1 \\
  6 & U1 ben 2 & U2 ben 2 \\
  7 & U2 ben 3 & U2 ben 4 \\
  8 & U1 ben 4 & U2 ben 5 \\
  9 & U1 ben 4 & U2 ben 10 \\
  10 & strömbrytare C & U2 ben 9 \\
  11 & U2 ben 6 & U3 ben 1 \\
  12 & U2 ben 8 & U3 ben 2 \\
  13 & U3 ben 3 & lysdiod F \\
  14 & U1 ben 5, 9, 11, 13 & GND \\
  15 & U2 ben 12, 13 & GND \\
  16 & U3 ben 4, 5, 9, 10, 12, 13 & GND \\
\end{narrowtable}

\noindent\code{B'} från U1 ben~4 används av två grindar, så två rader utgår från samma ben. På
kopplingsdäcket sätter du båda sladdarna i kolumnen för U1 ben~4.

\pt{c}Fläkten ska starta på hög temperatur när \code{T = 1}, \code{W = 0} och \code{B = 0}. Om
\code{C = 1} samtidigt startar den ändå, tack vare \code{CB'}. Den felaktiga raden är alltså
\textbf{\code{TWCB = 1000}}, där bara termen \code{TW'B'} kan göra \code{F = 1}. Ställ in den och
mät längs den termen: U2 ben~3 (\code{TW'}, ska vara 1) och U2 ben~6 (\code{TW'B'}, ska vara 1),
och sedan U3 ben~1.

\answer{c5:ex:10}{Kontroll: en utgång under last}
\pt{a}Med 5~V på utgången och 2~V över lysdioden ligger 3~V över motståndet:
\[
  I = \frac{5 - 2}{330} \approx 9,1\ \text{mA}
\]

\pt{b}Typiska mätvärden på en station är ungefär: hög utgång \textbf{4,4--4,7~V}, lysdiod
\textbf{1,9--2,0~V}, låg utgång \textbf{0,0--0,1~V}. Dina egna värden är svaret; skriv upp dem.

\pt{c}Med till exempel 4,5~V på utgången och 1,95~V över lysdioden:
\[
  I = \frac{4,5 - 1,95}{330} \approx 7,7\ \text{mA}
\]

\pt{d}Antagandet att utgången ger \textbf{exakt 5~V} var fel. En hög utgång har en viss inre
resistans, och när den levererar ström tappar den en del av spänningen: med siffrorna ovan ungefär
0,5~V vid 7,7~mA, alltså runt 65~Ω. Lysdiodens spänning kan också skilja sig lite från 2~V,
beroende på färg och exemplar.

Den uppmätta höga nivån ligger väl över databladets $V_{\text{OH}}$ på 3,84~V vid 4~mA, trots att
strömmen här är dubbelt så stor. Den låga nivån ligger nästan på 0~V, långt under $V_{\text{OL}}$
på 0,33~V, eftersom lysdioden är släckt när utgången är låg och ingen ström går genom utgången
alls. Slutsatsen: skillnaden mellan kalkyl och mätning är några tiondels volt och ett par
milliampere, och den kommer från utgångens inre resistans. Den spelar ingen roll för lysdioden, men
den visar varför en utgång inte ska belastas mer än nödvändigt.

\answer{c5:ex:11}{Fel krets i sockeln}
\pt{a}Ben~1 på en 74HC02 är \textbf{utgången} på grind~1, \code{1Y}.

\pt{b}När NOR-grindens ingångar (ben 2 och 3) båda är 0 försöker grinden ge en etta på ben~1. Om
strömbrytaren samtidigt håller ben~1 på 0~V kortsluts utgången mot jord genom strömbrytaren.
Strömmen begränsas bara av utgångstransistorns inre resistans, några tiotal ohm, och blir tiotals
milliampere: över kretsens gränsvärde på 25~mA. Kretsen blir varm och kan förstöras. Samma sak
händer åt andra hållet när strömbrytaren ger 5~V och grinden försöker ge 0.

\pt{c}Ben~3 är en \textbf{ingång} på en 74HC02 (\code{1B}). En ingång driver ingenting, så
lysdioden lyser aldrig, vad strömbrytarna än står på. Lysdioden och motståndet håller dessutom
ben~3 nära jord, så NOR-grinden ser ingången som 0. Kombinationen, en lysdiod som aldrig lyser och
en krets som blir varm, är kännetecknet för en krets med fel benplacering.
